%% kapa: exam=1 round=31 year=2020 period=2 form=A subject=appraisal_law subject_ko=감정평가관계법규 seq=18 number=18
%% source: Q-Net 감정평가사 기출문제 (cst00309) artlSeq=5207219 file=extracted/1008/5207219/A형/A형/2020년도 제31회 감정평가사 1차 시험_2교시_A형.pdf
%% conversion: kordoc 4.24.51 -> tools/gen1.py (auto); review=proofread-sonnet; flags=none
\begin{question}{18}
부동산 가격공시에 관한 법령상 주택가격의 공시에 관한 설명으로 옳은 것은?
\begin{choices}
  \item 국토교통부장관은 표준주택을 선정할 때에는 일반적으로 유사하다고 인정되는 일단의 공동주택 중에서 해당 일단의 공동주택을 대표할 수 있는 주택을 선정하여야 한다.
  \item 국토교통부장관은 표준주택가격을 조사ㆍ산정하고자 할 때에는 한국감정원 또는 둘 이상의 감정평가업자에게 의뢰한다.
  \item 표준주택가격은 국가ㆍ지방자치단체 등이 과세업무와 관련하여 주택의 가격을 산정하는 경우에 그 기준으로 활용하여야 한다.
  \item 표준주택가격의 공시사항에는 지목, 도로 상황이 포함되어야 한다.
  \item 개별주택가격 결정ㆍ공시에 소요되는 비용은 75퍼센트 이내에서 지방자치단체가 보조할 수 있다.
\end{choices}
\end{question}
\answer{4}
